\chapter{Direito Administrativo}

\section{Estado, governo e Administração Pública}

\termo{Estado} é pessoa jurídica soberana organizada por povo, território e poder político. \termo{Governo} corresponde à direção política superior, formulação de prioridades e exercício de funções políticas. \termo{Administração Pública}, em sentido subjetivo, designa órgãos, agentes e entidades; em sentido objetivo, a própria atividade administrativa.

\begin{teoria}[chave de prova]
A banca costuma misturar \textbf{quem exerce} a atividade com \textbf{qual atividade} é exercida. ``Administração Pública'' pode significar estrutura (sentido subjetivo/orgânico) ou função administrativa (sentido objetivo/material).
\end{teoria}

\section{Direito Administrativo: conceito, objeto e fontes}

Direito Administrativo disciplina a função administrativa, a organização da Administração, seus agentes, bens, serviços, poderes, controles e relações com particulares. Fontes incluem Constituição, leis, regulamentos, jurisprudência, doutrina e princípios, com forças normativas distintas.

\section{Regime jurídico-administrativo}

A atuação administrativa combina \termo{prerrogativas} para proteger o interesse público e \termo{sujeições} destinadas a preservar legalidade, direitos fundamentais e controle.

\subsection{Princípios expressos}

O art. 37 da Constituição consagra legalidade, impessoalidade, moralidade, publicidade e eficiência. A legalidade administrativa é mais restritiva que a autonomia privada: o agente público atua dentro das competências e finalidades conferidas pelo ordenamento.

\subsection{Princípios implícitos e correlatos}

Razoabilidade, proporcionalidade, segurança jurídica, motivação, autotutela, continuidade, especialidade, proteção da confiança, contraditório e ampla defesa aparecem em diferentes contextos normativos e jurisprudenciais.

\section{Organização administrativa}

\subsection{Centralização e descentralização}

Centralização ocorre quando o próprio ente político executa a atividade por seus órgãos. Descentralização transfere execução/titularidade nos termos legais a outra pessoa jurídica ou particular.

\subsection{Concentração e desconcentração}

Desconcentração distribui competências internamente entre órgãos da mesma pessoa jurídica. Não cria nova pessoa. Concentração faz o movimento inverso.

\subsection{Administração direta e indireta}

Administração direta integra os entes políticos e seus órgãos. A indireta é composta, na classificação clássica, por autarquias, fundações públicas, empresas públicas e sociedades de economia mista, observadas suas peculiaridades de criação/autorização, regime e finalidade.

\begin{atencao}
Órgão não possui personalidade jurídica própria. Entidade possui. Essa diferença é recorrente em questões de desconcentração e descentralização.
\end{atencao}

\section{Atos administrativos}

\subsection{Elementos/requisitos}

A abordagem clássica trabalha com competência, finalidade, forma, motivo e objeto. A validade depende do atendimento ao regime jurídico aplicável.

\subsection{Atributos}

Presunção de legitimidade/veracidade, imperatividade, autoexecutoriedade e tipicidade aparecem em doutrina, mas \textbf{nem todo ato possui todos os atributos}. Autoexecutoriedade, por exemplo, depende de previsão legal ou situação que a justifique.

\subsection{Mérito administrativo}

Mérito relaciona-se a conveniência e oportunidade nos espaços legítimos de discricionariedade. Discricionariedade não significa liberdade fora da lei; continua sujeita a competência, finalidade, motivos, proporcionalidade e controle.

\subsection{Anulação, revogação e convalidação}

\termo{Anulação} decorre de ilegalidade e pode ser feita pela Administração, sujeita ao regime jurídico e à segurança jurídica, ou pelo Judiciário quando provocado. \termo{Revogação} atinge ato válido por conveniência/oportunidade e é ato da Administração, não do Judiciário no exercício típico do controle de legalidade.

Convalidação sana vícios passíveis de correção, quando não houver lesão ao interesse público ou prejuízo a terceiros, nos termos legais.

\subsection{Decadência administrativa}

A Lei nº 9.784/1999 estabelece prazo decadencial para anulação, pela Administração federal, de atos favoráveis ao destinatário, ressalvada má-fé, observadas as regras legais. Em prova, diferencie decadência administrativa de prescrição de pretensões.

\section{Agentes públicos}

Agentes públicos abrangem múltiplas categorias: agentes políticos, servidores estatutários, empregados públicos, temporários e particulares em colaboração, conforme classificação doutrinária.

\subsection{Cargo, emprego e função}

Cargo público integra estrutura estatutária. Emprego público é regido, em regra, pela CLT com incidência de normas de direito público. Função pública é conjunto de atribuições, podendo existir sem cargo em hipóteses específicas.

\subsection{Provimento e vacância}

A Lei nº 8.112/1990 disciplina, no âmbito federal, formas de provimento e vacância, direitos, deveres e processo disciplinar. Nomeação é forma originária de provimento; outras formas são derivadas conforme previsão legal.

\subsection{Efetividade, estabilidade e vitaliciedade}

Efetividade decorre da forma de investidura em cargo efetivo. Estabilidade é garantia constitucional adquirida sob requisitos próprios. Vitaliciedade possui regime especial para determinados agentes. Não são sinônimos.

\subsection{Responsabilidade do servidor}

Responsabilidades civil, penal e administrativa podem coexistir e possuem relativa independência, ressalvados efeitos legais de decisões em outras esferas.

\section{Poderes administrativos}

\subsection{Hierárquico}

Organiza internamente competências, delegação, avocação, fiscalização e revisão dentro dos limites legais. Não existe hierarquia entre pessoas jurídicas apenas por vínculo de supervisão finalística.

\subsection{Disciplinar}

Permite apurar infrações e aplicar sanções a agentes e outros sujeitos submetidos a vínculo jurídico especial com a Administração.

\subsection{Regulamentar/normativo}

Permite editar atos gerais para fiel execução da lei e, nas hipóteses constitucionais, decretos autônomos dentro de limites específicos.

\subsection{Polícia}

Poder de polícia condiciona ou restringe liberdade/propriedade em benefício de interesses coletivos juridicamente protegidos. Seus ciclos, delegabilidade e atributos variam conforme atividade e entendimento jurisprudencial.

\subsection{Abuso de poder}

Excesso de poder ocorre quando agente ultrapassa competência; desvio de finalidade ocorre quando usa competência para fim diverso do previsto.

\section{Responsabilidade civil do Estado}

A Constituição adota responsabilidade objetiva das pessoas jurídicas de direito público e das pessoas jurídicas de direito privado prestadoras de serviços públicos pelos danos que seus agentes, nessa qualidade, causarem a terceiros, assegurado direito de regresso contra o responsável nos casos previstos.

\subsection{Elementos}

Conduta administrativa, dano e nexo causal são essenciais. Causas excludentes/atenuantes podem romper ou reduzir o nexo conforme o caso. A responsabilidade por omissão possui tratamento jurisprudencial mais complexo e não deve ser reduzida a fórmula única.

\section{Serviços públicos}

Serviço público é atividade destinada a satisfazer necessidade coletiva sob regime jurídico próprio. Princípios recorrentes incluem continuidade, generalidade, modicidade, eficiência, atualidade e segurança, conforme regime.

Concessão e permissão de serviços públicos possuem disciplina legal específica e pressupõem licitação nas hipóteses legais.

\section{Controle da Administração}

\subsection{Controle administrativo}

Autotutela permite à Administração controlar seus próprios atos, anulando ilegais e revogando válidos inconvenientes/oportunos dentro dos limites jurídicos.

\subsection{Controle judicial}

Judiciário controla legalidade e juridicidade quando provocado, sem substituir legitimamente a Administração no mérito puro de conveniência e oportunidade.

\subsection{Controle legislativo}

O Poder Legislativo exerce controle político e financeiro nos termos constitucionais, com auxílio dos Tribunais de Contas no controle externo.

\section{Improbidade administrativa}

A Lei nº 8.429/1992, profundamente alterada pela Lei nº 14.230/2021, disciplina atos de improbidade, sanções, procedimento e prescrição. A redação atual exige atenção especial ao elemento subjetivo e à tipificação legal.

\begin{cebraspe}
Material antigo que trata culpa como suficiente de forma genérica para improbidade pode estar desatualizado. A redação pós-Lei nº 14.230/2021 exige estudo do dolo nos tipos legais e da jurisprudência correspondente.
\end{cebraspe}

\section{Processo administrativo federal — Lei nº 9.784/1999}

A lei estabelece princípios, direitos, deveres, competência, impedimento/suspeição, forma, motivação, instrução, decisão, recursos, anulação, revogação e convalidação no âmbito federal.

\subsection{Motivação}

Motivação explicita fundamentos de fato e de direito. É especialmente relevante em atos que afetem direitos, imponham sanções, decidam recursos ou se afastem de entendimentos anteriores, conforme disciplina legal.

\subsection{Delegação e avocação}

Delegação é admitida dentro das condições legais, inclusive entre órgãos sem relação hierárquica em hipóteses previstas. Certas matérias não podem ser delegadas. Avocação é medida excepcional e temporária, sujeita a justificativa e limites.

\section{Licitações e contratos — Lei nº 14.133/2021}

A nova Lei de Licitações estrutura planejamento, seleção, contratação, execução e responsabilização.

\subsection{Princípios e planejamento}

Planejamento, segregação de funções, transparência, motivação, julgamento objetivo, competitividade, economicidade e eficiência devem orientar a contratação.

\subsection{Fases da licitação}

O procedimento segue fases definidas na lei, com possibilidade de inversões nos casos legalmente admitidos. O candidato deve entender a lógica, não decorar sequência sem contexto.

\subsection{Contratação direta}

Dispensa e inexigibilidade são hipóteses diferentes. Inexigibilidade pressupõe inviabilidade de competição; dispensa ocorre em hipóteses legais mesmo quando competição poderia ser possível.

\subsection{Execução contratual}

Gestão e fiscalização verificam cumprimento do objeto, medições, alterações, sanções e recebimento. Alterações e reequilíbrio dependem de fundamento e requisitos próprios.

\section{Mapa mental funcional}

\mapafuncional{Direito Administrativo}
{Administração atua com prerrogativas, mas sempre sob competência, finalidade, controle e direitos}
{Estrutura: Estado $\rightarrow$ Administração direta/indireta $\rightarrow$ órgãos/entidades.\\Atuação: ato $\rightarrow$ poder $\rightarrow$ agente $\rightarrow$ serviço/contrato.\\Controle: autotutela + Judiciário + Legislativo/TCE.\\Responsabilização: Estado + servidor + improbidade.}
{Ato ilegal: pense anulação.\\Ato válido inconveniente: revogação.\\Desconcentração: mesma pessoa; descentralização: outra pessoa.\\Poder exercido fora da competência: excesso; fim impróprio: desvio.\\Contratação sem competição: diferencie dispensa de inexigibilidade.}
{Discricionariedade $\neq$ arbitrariedade.\\Efetividade $\neq$ estabilidade.\\Órgão $\neq$ entidade.\\Revogação $\neq$ anulação.\\Supervisão finalística $\neq$ hierarquia.}
{Qual é a competência?\\O ato é válido?\\Há direito adquirido/segurança jurídica?\\A pessoa é órgão ou entidade?\\O controle é de legalidade ou mérito?\\A lei vigente pós-2021 foi considerada?}

\begin{resumo}
\textbf{Regra de prova:} quase todas as pegadinhas de Direito Administrativo nascem de confundir categorias próximas: anulação/revogação, descentralização/desconcentração, cargo/emprego, poder hierárquico/supervisão, dispensa/inexigibilidade.
\end{resumo}

\section{Banco de questões — Direito Administrativo}

\subsection*{N1 — Fundamento competitivo}

\begin{questao}{\nivel{N1} 04.01 — Cebraspe/Certo ou Errado}
Em sentido subjetivo, Administração Pública designa os órgãos, agentes e entidades que exercem função administrativa; em sentido objetivo, refere-se à própria atividade administrativa.
\end{questao}

\begin{questao}{\nivel{N1} 04.02 — Cebraspe/Certo ou Errado}
A criação de novos órgãos dentro de uma mesma pessoa jurídica caracteriza descentralização administrativa, pois distribui competências entre centros distintos.
\end{questao}

\begin{questao}{\nivel{N1} 04.03 — Cebraspe/Certo ou Errado}
Órgãos públicos, embora possuam competências próprias, não dispõem de personalidade jurídica autônoma em relação à pessoa jurídica que integram.
\end{questao}

\begin{questao}{\nivel{N1} 04.04 — Cebraspe/Certo ou Errado}
Todo ato administrativo possui, necessariamente, imperatividade e autoexecutoriedade.
\end{questao}

\begin{questao}{\nivel{N1} 04.05 — Cebraspe/Certo ou Errado}
A Administração pode anular ato ilegal e revogar ato válido por razões de conveniência e oportunidade, observados os limites jurídicos aplicáveis.
\end{questao}

\begin{questao}{\nivel{N1} 04.06 — Cebraspe/Certo ou Errado}
O Poder Judiciário, no exercício típico do controle jurisdicional, pode revogar ato administrativo válido sempre que considerar inconveniente a escolha feita pela Administração.
\end{questao}

\begin{questao}{\nivel{N1} 04.07 — Cebraspe/Certo ou Errado}
Efetividade e estabilidade são conceitos equivalentes: todo servidor ocupante de cargo efetivo é estável desde a posse.
\end{questao}

\begin{questao}{\nivel{N1} 04.08 — Cebraspe/Certo ou Errado}
A supervisão finalística exercida sobre entidade da Administração indireta não se confunde com hierarquia existente entre órgãos da mesma pessoa jurídica.
\end{questao}

\begin{questao}{\nivel{N1} 04.09 — Cebraspe/Certo ou Errado}
Na inexigibilidade de licitação, a característica central é a inviabilidade de competição; na dispensa, a lei admite contratação direta em hipóteses previstas mesmo quando competição poderia existir.
\end{questao}

\begin{questao}{\nivel{N1} 04.10 — Cebraspe/Certo ou Errado}
Após as alterações promovidas pela Lei nº 14.230/2021, a mera culpa do agente continua suficiente, em qualquer hipótese, para caracterizar ato de improbidade administrativa.
\end{questao}

\subsection*{N2 — Aplicação}

\begin{questao}{\nivel{N2} 04.11 — FGV/organização administrativa}
Um ministério cria internamente uma secretaria especializada, sem instituição de nova pessoa jurídica. O fenômeno é:
\begin{enumerate}[label=\Alph*)]
\item descentralização por serviços;
\item desconcentração administrativa;
\item privatização;
\item delegação contratual de serviço público;
\item criação de autarquia.
\end{enumerate}
\end{questao}

\begin{questao}{\nivel{N2} 04.12 — FCC/administração indireta}
Assinale a alternativa correta.
\begin{enumerate}[label=\Alph*)]
\item Autarquia é órgão sem personalidade jurídica.
\item Empresa pública possui necessariamente capital misto público-privado.
\item Sociedade de economia mista integra a Administração indireta e possui personalidade jurídica própria.
\item Fundação pública nunca integra a Administração indireta.
\item Toda entidade indireta está subordinada hierarquicamente ao ministério supervisor.
\end{enumerate}
\end{questao}

\begin{questao}{\nivel{N2} 04.13 — Cebraspe/Certo ou Errado}
A discricionariedade administrativa existe apenas dentro dos espaços conferidos pelo ordenamento e continua sujeita a controle de finalidade, motivos, proporcionalidade e juridicidade.
\end{questao}

\begin{questao}{\nivel{N2} 04.14 — FGV/atos administrativos}
A autoridade identifica vício de competência em ato praticado por agente incompetente, mas a competência não era exclusiva e não houve prejuízo ao interesse público nem a terceiros. Em tese, a providência juridicamente possível é:
\begin{enumerate}[label=\Alph*)]
\item revogação obrigatória;
\item convalidação, se preenchidos os requisitos legais;
\item prescrição da competência;
\item conversão automática em lei;
\item declaração de inexistência do ato em qualquer caso.
\end{enumerate}
\end{questao}

\begin{questao}{\nivel{N2} 04.15 — FCC/anulação e revogação}
Um ato válido deixou de ser conveniente para a Administração, sem gerar direito adquirido à sua manutenção. A medida típica é:
\begin{enumerate}[label=\Alph*)]
\item anulação por ilegalidade;
\item revogação por mérito administrativo;
\item convalidação;
\item cassação judicial obrigatória;
\item declaração de nulidade absoluta pelo particular.
\end{enumerate}
\end{questao}

\begin{questao}{\nivel{N2} 04.16 — Cebraspe/Certo ou Errado}
A presunção de legitimidade dos atos administrativos é relativa e admite prova em contrário.
\end{questao}

\begin{questao}{\nivel{N2} 04.17 — FGV/poderes administrativos}
Um agente competente para aplicar determinada sanção utiliza essa atribuição para retaliar servidor por divergência pessoal. O vício predominante caracteriza:
\begin{enumerate}[label=\Alph*)]
\item excesso de poder;
\item desvio de finalidade;
\item delegação válida;
\item discricionariedade legítima;
\item convalidação tácita.
\end{enumerate}
\end{questao}

\begin{questao}{\nivel{N2} 04.18 — FCC/poder hierárquico}
O poder hierárquico permite, observados os limites legais:
\begin{enumerate}[label=\Alph*)]
\item criar hierarquia entre autarquia e ministério supervisor;
\item distribuir funções, fiscalizar subordinados e rever atos internos;
\item aplicar sanção penal a qualquer cidadão;
\item editar lei formal;
\item afastar controle judicial.
\end{enumerate}
\end{questao}

\begin{questao}{\nivel{N2} 04.19 — Cebraspe/Certo ou Errado}
Excesso de poder relaciona-se ao ultrapassamento da competência; desvio de finalidade, ao uso da competência para finalidade diversa da prevista juridicamente.
\end{questao}

\begin{questao}{\nivel{N2} 04.20 — FGV/responsabilidade civil}
Veículo oficial, conduzido por servidor no exercício da função, causa dano a terceiro por atuação do agente. À luz do art. 37, §6º, da Constituição, a responsabilidade da pessoa jurídica é, em regra:
\begin{enumerate}[label=\Alph*)]
\item subjetiva, exigindo prova de dolo do Estado;
\item objetiva, sem prejuízo do direito de regresso nos casos cabíveis;
\item inexistente, pois o agente responde sozinho;
\item penal;
\item exclusivamente contratual.
\end{enumerate}
\end{questao}

\begin{questao}{\nivel{N2} 04.21 — Cebraspe/Certo ou Errado}
A responsabilidade objetiva do Estado dispensa a demonstração de nexo causal entre a atuação estatal e o dano.
\end{questao}

\begin{questao}{\nivel{N2} 04.22 — FCC/agentes públicos}
A estabilidade constitucional do servidor ocupante de cargo efetivo:
\begin{enumerate}[label=\Alph*)]
\item decorre imediatamente da nomeação;
\item confunde-se com efetividade;
\item depende do atendimento aos requisitos constitucionais próprios e não nasce automaticamente com a posse;
\item alcança todo ocupante de cargo em comissão;
\item é idêntica à vitaliciedade.
\end{enumerate}
\end{questao}

\begin{questao}{\nivel{N2} 04.23 — FGV/Lei 9.784}
A autoridade pretende avocar competência de órgão inferior por tempo indeterminado e sem justificativa. A medida é:
\begin{enumerate}[label=\Alph*)]
\item compatível com a lei, pois avocação é livre;
\item incompatível, pois a avocação é excepcional, temporária e deve ser justificada nos termos legais;
\item obrigatória quando houver hierarquia;
\item equivalente à delegação;
\item possível apenas por decisão judicial.
\end{enumerate}
\end{questao}

\begin{questao}{\nivel{N2} 04.24 — Cebraspe/Certo ou Errado}
Na Lei nº 9.784/1999, determinadas matérias, como decisão de recursos administrativos e edição de atos de caráter normativo, não podem ser objeto de delegação.
\end{questao}

\begin{questao}{\nivel{N2} 04.25 — FCC/licitações}
Na Lei nº 14.133/2021, contratação direta por inviabilidade de competição enquadra-se como:
\begin{enumerate}[label=\Alph*)]
\item dispensa;
\item inexigibilidade;
\item pregão obrigatório;
\item concorrência simplificada;
\item diálogo competitivo necessariamente.
\end{enumerate}
\end{questao}

\subsection*{N3 — Integração de concurso}

\begin{questao}{\nivel{N3} 04.26 — FGV/assertivas sobre atos}
Considere:

I. A anulação decorre de ilegalidade.

II. A revogação pressupõe ato válido e juízo de conveniência/oportunidade.

III. O Judiciário pode revogar ato administrativo válido no controle típico de legalidade.

Assinale a alternativa correta.
\begin{enumerate}[label=\Alph*)]
\item Apenas I.
\item Apenas II.
\item Apenas I e II.
\item Apenas II e III.
\item I, II e III.
\end{enumerate}
\end{questao}

\begin{questao}{\nivel{N3} 04.27 — Cebraspe/Certo ou Errado}
A decadência administrativa prevista na Lei nº 9.784/1999 pode limitar o poder de anular atos favoráveis ao destinatário, ressalvada a hipótese de má-fé e demais condições legais.
\end{questao}

\begin{questao}{\nivel{N3} 04.28 — FCC/motivação}
A Administração aplica sanção a servidor por decisão que apenas afirma ``por interesse público'', sem indicar fatos ou fundamentos. A deficiência principal é de:
\begin{enumerate}[label=\Alph*)]
\item competência legislativa;
\item motivação;
\item personalidade jurídica;
\item publicidade da lei;
\item descentralização.
\end{enumerate}
\end{questao}

\begin{questao}{\nivel{N3} 04.29 — FGV/autotutela}
A Administração descobre ato ilegal próprio e, sem depender de provocação judicial, inicia procedimento para anulá-lo, respeitando contraditório quando necessário. A atuação decorre do princípio da:
\begin{enumerate}[label=\Alph*)]
\item continuidade;
\item autotutela;
\item especialidade;
\item supremacia do Legislativo;
\item indisponibilidade do processo judicial.
\end{enumerate}
\end{questao}

\begin{questao}{\nivel{N3} 04.30 — Cebraspe/Certo ou Errado}
A autotutela administrativa não torna dispensáveis contraditório e ampla defesa quando a invalidação afeta situação jurídica individual consolidada em contexto que exija tais garantias.
\end{questao}

\begin{questao}{\nivel{N3} 04.31 — FCC/serviço público}
A continuidade do serviço público:
\begin{enumerate}[label=\Alph*)]
\item significa impossibilidade absoluta de qualquer interrupção;
\item é princípio relevante, mas admite hipóteses de interrupção previstas no regime jurídico aplicável;
\item elimina o direito de greve em qualquer situação sem mediação constitucional;
\item impede manutenção programada;
\item dispensa planejamento de contingência.
\end{enumerate}
\end{questao}

\begin{questao}{\nivel{N3} 04.32 — FGV/responsabilidade por omissão}
Sobre responsabilidade estatal por omissão, assinale a melhor afirmação.
\begin{enumerate}[label=\Alph*)]
\item Toda omissão gera responsabilidade objetiva automática.
\item O tema exige análise do dever jurídico de agir, do nexo e do regime jurisprudencial aplicável, não cabendo fórmula única para qualquer omissão.
\item O Estado jamais responde por omissão.
\item Só há responsabilidade se houver contrato.
\item O agente responde sozinho.
\end{enumerate}
\end{questao}

\begin{questao}{\nivel{N3} 04.33 — Cebraspe/Certo ou Errado}
A existência de causa exclusiva da vítima pode, conforme o caso, romper o nexo causal e afastar a responsabilidade estatal.
\end{questao}

\begin{questao}{\nivel{N3} 04.34 — FCC/improbidade}
Após a Lei nº 14.230/2021, a análise de improbidade administrativa deve, entre outros aspectos:
\begin{enumerate}[label=\Alph*)]
\item presumir dolo de qualquer ilegalidade;
\item verificar tipificação legal e elemento subjetivo exigido, não equiparando toda ilegalidade a improbidade;
\item aplicar automaticamente a Lei Penal;
\item ignorar prescrição;
\item tratar culpa simples como suficiente em qualquer tipo.
\end{enumerate}
\end{questao}

\begin{questao}{\nivel{N3} 04.35 — FGV/licitações}
Um órgão escolhe fornecedor específico por preferência técnica interna, embora existam vários capazes de competir, e tenta enquadrar a contratação como inexigível. O problema central é:
\begin{enumerate}[label=\Alph*)]
\item inexigibilidade exige inviabilidade de competição demonstrada, não mera preferência;
\item todo fornecedor conhecido gera inexigibilidade;
\item competição é vedada em TI;
\item escolha técnica elimina necessidade de motivação;
\item inexigibilidade e dispensa são sinônimos.
\end{enumerate}
\end{questao}

\begin{questao}{\nivel{N3} 04.36 — Cebraspe/Certo ou Errado}
Na contratação direta, a ausência de licitação não elimina o dever de instrução, motivação, justificativa de preço e demonstração do enquadramento legal.
\end{questao}

\begin{questao}{\nivel{N3} 04.37 — FCC/controle}
O Tribunal de Contas, no exercício do controle externo, integra:
\begin{enumerate}[label=\Alph*)]
\item a Administração direta do Poder Executivo;
\item o sistema constitucional de controle externo exercido pelo Legislativo com seu auxílio;
\item o Poder Judiciário;
\item a função de autotutela do órgão auditado;
\item exclusivamente o controle interno.
\end{enumerate}
\end{questao}

\begin{questao}{\nivel{N3} 04.38 — FGV/supervisão finalística}
Ministério determina diretamente a um presidente de autarquia como decidir processo técnico específico, invocando ``poder hierárquico''. A justificativa é problemática porque:
\begin{enumerate}[label=\Alph*)]
\item autarquia não integra o Estado;
\item não há, em regra, subordinação hierárquica entre pessoas jurídicas distintas; existe vinculação/supervisão nos termos legais;
\item autarquias não estão sujeitas a controle algum;
\item ministérios não podem exercer supervisão;
\item toda autarquia é independente como Poder.
\end{enumerate}
\end{questao}

\begin{questao}{\nivel{N3} 04.39 — Cebraspe/Certo ou Errado]
A discricionariedade pode envolver escolha legítima entre alternativas legais, mas a motivação declarada pode ser controlada quanto à veracidade dos pressupostos fáticos que a sustentam.
\end{questao}

\begin{questao}{\nivel{N3} 04.40 — FCC/Lei 8.112}
A responsabilidade administrativa de servidor:
\begin{enumerate}[label=\Alph*)]
\item exclui sempre a responsabilidade civil e penal;
\item pode coexistir com responsabilidades civil e penal, observadas as relações legais entre as instâncias;
\item somente existe se houver condenação penal;
\item nunca depende de processo disciplinar;
\item é objetiva em qualquer infração.
\end{enumerate}
\end{questao}

\subsection*{N4 — Auditoria / avançado}

\begin{questao}{\nivel{N4} 04.41 — Cebraspe/caso integrado}
Um ato favorável foi praticado há anos com vício, sem indício de má-fé do beneficiário. Julgue o item:

Antes de simplesmente anulá-lo, a Administração deve considerar o regime decadencial e a segurança jurídica aplicáveis, além do contraditório quando pertinente.
\end{questao}

\begin{questao}{\nivel{N4} 04.42 — FGV/controle de mérito}
O Judiciário reconhece que determinado ato discricionário é legal, mas entende que outra escolha seria administrativamente ``mais eficiente''. A providência constitucionalmente mais adequada é:
\begin{enumerate}[label=\Alph*)]
\item substituir a escolha por aquela que o juiz considera melhor;
\item respeitar o espaço legítimo de mérito, salvo vício de juridicidade controlável;
\item revogar o ato;
\item assumir a gestão do órgão;
\item determinar contratação específica sem base legal.
\end{enumerate}
\end{questao}

\begin{questao}{\nivel{N4} 04.43 — FCC/contratação direta}
Auditoria identifica contratação por inexigibilidade com três fornecedores aptos e ausência de demonstração de exclusividade ou inviabilidade de competição. O achado mais adequado é:
\begin{enumerate}[label=\Alph*)]
\item ``Toda inexigibilidade é ilegal'';
\item ``O processo não demonstra a inviabilidade de competição necessária ao enquadramento e requer revisão da motivação, pesquisa de mercado e responsabilização conforme o caso'';
\item ``Basta reduzir o preço'';
\item ``A contratação é válida se o fornecedor for conhecido'';
\item ``Deve-se trocar a modalidade para dispensa sem nova análise''.
\end{enumerate}
\end{questao}

\begin{questao}{\nivel{N4} 04.44 — Cebraspe/Certo ou Errado}
Uma ilegalidade administrativa pode existir sem que, automaticamente, estejam preenchidos todos os requisitos de um ato de improbidade após a reforma da Lei nº 8.429/1992.
\end{questao}

\begin{questao}{\nivel{N4} 04.45 — FGV/poder de polícia}
Órgão aplica restrição a atividade privada sem competência legal e invoca genericamente ``poder de polícia''. O vício principal é:
\begin{enumerate}[label=\Alph*)]
\item poder de polícia afasta legalidade;
\item falta de competência, pois prerrogativa administrativa depende de fundamento jurídico;
\item ausência de personalidade do particular;
\item inexistência de interesse público em qualquer fiscalização;
\item excesso de publicidade.
\end{enumerate}
\end{questao}

\begin{questao}{\nivel{N4} 04.46 — Cebraspe/Certo ou Errado}
A motivação de ato discricionário pode vincular a Administração aos motivos declarados, de modo que a inexistência ou falsidade desses pressupostos seja relevante para o controle de validade.
\end{questao}

\begin{questao}{\nivel{N4} 04.47 — FCC/responsabilidade estatal}
Um terceiro sofre dano decorrente de atuação de empregado de concessionária de serviço público, nessa qualidade. Em tese, o regime constitucional do art. 37, §6º, alcança:
\begin{enumerate}[label=\Alph*)]
\item apenas pessoas jurídicas de direito público;
\item também pessoas jurídicas de direito privado prestadoras de serviços públicos, nos termos constitucionais;
\item exclusivamente o agente pessoa física;
\item somente contratos de obra;
\item nenhuma concessionária.
\end{enumerate}
\end{questao}

\begin{questao}{\nivel{N4} 04.48 — FGV/processo administrativo}
Autoridade indefere recurso administrativo com texto padronizado que não enfrenta nenhum argumento novo e relevante apresentado pelo recorrente. O principal ponto de controle é:
\begin{enumerate}[label=\Alph*)]
\item descentralização;
\item suficiência da motivação e observância do devido processo administrativo;
\item personalidade jurídica do órgão;
\item licitação;
\item responsabilidade objetiva.
\end{enumerate}
\end{questao}

\begin{questao}{\nivel{N4} 04.49 — Cebraspe/Certo ou Errado}
O princípio da eficiência autoriza a Administração a afastar exigência legal sempre que a solução informal for mais rápida e barata.
\end{questao}

\begin{questao}{\nivel{N4} 04.50 — FGV/matriz de achado}
Auditoria identifica renovação contratual por inércia, sem análise de desempenho, vantajosidade ou alternativas, embora o serviço seja crítico. O achado mais adequado é:
\begin{enumerate}[label=\Alph*)]
\item ``Prorrogar contratos é proibido'';
\item ``A ausência de motivação e avaliação de vantajosidade/desempenho fragiliza a decisão de prorrogação e pode perpetuar contratação desvantajosa; recomenda-se processo decisório documentado e análise de alternativas'';
\item ``Deve-se trocar de fornecedor obrigatoriamente'';
\item ``Basta exigir desconto de 1\%'';
\item ``Serviço crítico deve ser sempre internalizado''.
\end{enumerate}
\end{questao}


\section{Treino discursivo}
\begin{discursiva}[ato e controle]
Diferencie anulação, revogação e convalidação de atos administrativos, indicando fundamento, autoridade competente, efeitos e limites decorrentes da segurança jurídica.
\end{discursiva}

\section{Gabarito e resoluções comentadas — Direito Administrativo}

\begin{solucao}{04.01}\textbf{CERTO.} O sentido subjetivo/orgânico aponta os sujeitos da Administração; o objetivo/material, a atividade administrativa.\end{solucao}
\begin{solucao}{04.02}\textbf{ERRADO.} Distribuir competências dentro da mesma pessoa jurídica é desconcentração. Descentralização envolve outra pessoa jurídica ou particular.\end{solucao}
\begin{solucao}{04.03}\textbf{CERTO.} Órgão é centro de competências sem personalidade jurídica própria.\end{solucao}
\begin{solucao}{04.04}\textbf{ERRADO.} Nem todo ato é imperativo ou autoexecutório; esses atributos dependem da natureza e do regime do ato.\end{solucao}
\begin{solucao}{04.05}\textbf{CERTO.} Anulação trata ilegalidade; revogação, mérito de ato válido, respeitados direitos e limites.\end{solucao}
\begin{solucao}{04.06}\textbf{ERRADO.} O Judiciário controla juridicidade, não substitui a Administração no mérito legítimo para revogar ato válido.\end{solucao}
\begin{solucao}{04.07}\textbf{ERRADO.} Efetividade decorre da investidura em cargo efetivo; estabilidade é garantia adquirida sob requisitos constitucionais próprios.\end{solucao}
\begin{solucao}{04.08}\textbf{CERTO.} Vinculação e supervisão finalística não criam subordinação hierárquica entre pessoas jurídicas distintas.\end{solucao}
\begin{solucao}{04.09}\textbf{CERTO.} Inexigibilidade pressupõe inviabilidade de competição; dispensa decorre de hipótese legal de contratação direta.\end{solucao}
\begin{solucao}{04.10}\textbf{ERRADO.} A reforma de 2021 reforçou o dolo nos tipos de improbidade, afastando a tese de culpa suficiente em qualquer hipótese.\end{solucao}

\begin{solucao}{04.11}\textbf{B.} Criar secretaria interna é desconcentrar competências sem criar nova pessoa jurídica.\end{solucao}
\begin{solucao}{04.12}\textbf{C.} Sociedade de economia mista é entidade com personalidade própria da Administração indireta. Autarquia também é entidade, não órgão.\end{solucao}
\begin{solucao}{04.13}\textbf{CERTO.} Discricionariedade é liberdade juridicamente delimitada, não imunidade ao controle.\end{solucao}
\begin{solucao}{04.14}\textbf{B.} Vício de competência não exclusiva pode ser convalidável, desde que presentes os requisitos legais e ausentes prejuízos impeditivos.\end{solucao}
\begin{solucao}{04.15}\textbf{B.} Ato válido que se torna inconveniente pode ser revogado pela Administração, observados limites jurídicos.\end{solucao}
\begin{solucao}{04.16}\textbf{CERTO.} A presunção é juris tantum, podendo ser afastada por prova.\end{solucao}
\begin{solucao}{04.17}\textbf{B.} O agente tem competência, mas a usa para finalidade imprópria: desvio de finalidade.\end{solucao}
\begin{solucao}{04.18}\textbf{B.} O poder hierárquico organiza e controla internamente relações de subordinação.\end{solucao}
\begin{solucao}{04.19}\textbf{CERTO.} Essa é a distinção clássica entre as modalidades de abuso de poder.\end{solucao}
\begin{solucao}{04.20}\textbf{B.} O art. 37, §6º, adota responsabilidade objetiva da pessoa jurídica, preservado o regresso contra o agente nos casos constitucionais.\end{solucao}
\begin{solucao}{04.21}\textbf{ERRADO.} Responsabilidade objetiva dispensa culpa do Estado, não dano e nexo causal.\end{solucao}
\begin{solucao}{04.22}\textbf{C.} Estabilidade não nasce da posse e não se confunde com efetividade ou vitaliciedade.\end{solucao}
\begin{solucao}{04.23}\textbf{B.} A avocação é excepcional e temporária, devendo ser motivada dentro dos limites legais.\end{solucao}
\begin{solucao}{04.24}\textbf{CERTO.} A Lei nº 9.784/1999 contém hipóteses de competência indelegável, entre elas matérias normativas e decisão de recursos.\end{solucao}
\begin{solucao}{04.25}\textbf{B.} Inviabilidade de competição é o núcleo da inexigibilidade.\end{solucao}

\begin{solucao}{04.26}\textbf{C.} I e II estão corretas. III é falsa porque o Judiciário não revoga ato válido por juízo de conveniência.\end{solucao}
\begin{solucao}{04.27}\textbf{CERTO.} O poder de autotutela convive com decadência e segurança jurídica, ressalvada má-fé nos termos legais.\end{solucao}
\begin{solucao}{04.28}\textbf{B.} A decisão sancionatória precisa expor fundamentos de fato e de direito suficientes.\end{solucao}
\begin{solucao}{04.29}\textbf{B.} Autotutela permite à Administração invalidar seus atos ilegais e rever atos válidos dentro do regime aplicável.\end{solucao}
\begin{solucao}{04.30}\textbf{CERTO.} A invalidação de situação individual pode exigir contraditório e ampla defesa, conforme o contexto e o regime jurídico.\end{solucao}
\begin{solucao}{04.31}\textbf{B.} Continuidade é princípio, mas o ordenamento admite interrupções justificadas e reguladas.\end{solucao}
\begin{solucao}{04.32}\textbf{B.} Omissões exigem análise do dever de agir e do nexo, com atenção ao tratamento jurisprudencial específico.\end{solucao}
\begin{solucao}{04.33}\textbf{CERTO.} Culpa exclusiva da vítima pode romper o nexo; culpa concorrente pode influenciar o resultado conforme o caso.\end{solucao}
\begin{solucao}{04.34}\textbf{B.} Nem toda ilegalidade é improbidade; é necessário verificar o tipo e o elemento subjetivo exigido na lei vigente.\end{solucao}
\begin{solucao}{04.35}\textbf{A.} Preferência administrativa não cria inviabilidade de competição. O processo deve demonstrar objetivamente o enquadramento.\end{solucao}
\begin{solucao}{04.36}\textbf{CERTO.} Contratação direta continua sujeita a instrução e motivação, inclusive quanto a preço e hipótese legal.\end{solucao}
\begin{solucao}{04.37}\textbf{B.} Tribunais de Contas auxiliam o Poder Legislativo no controle externo, conforme desenho constitucional.\end{solucao}
\begin{solucao}{04.38}\textbf{B.} Autarquia é pessoa jurídica própria; há vinculação e supervisão, não relação hierárquica comum com o ministério.\end{solucao}
\begin{solucao}{04.39}\textbf{CERTO.} A teoria dos motivos determinantes torna relevante a correspondência entre os motivos declarados e a realidade, inclusive em atos discricionários.\end{solucao}
\begin{solucao}{04.40}\textbf{B.} As responsabilidades podem coexistir, com independência relativa e efeitos legais entre instâncias.\end{solucao}

\begin{solucao}{04.41}\textbf{CERTO.} O poder de invalidar não é ilimitado no tempo e deve respeitar decadência, boa-fé e garantias processuais aplicáveis.\end{solucao}
\begin{solucao}{04.42}\textbf{B.} Eficiência não autoriza o Judiciário a substituir opção administrativa legítima por preferência própria sem vício jurídico.\end{solucao}
\begin{solucao}{04.43}\textbf{B.} O achado deve demonstrar a ausência do pressuposto da inexigibilidade e indicar análise da instrução e eventual responsabilização, sem generalizar.\end{solucao}
\begin{solucao}{04.44}\textbf{CERTO.} Ilegalidade e improbidade são categorias distintas; a segunda exige requisitos específicos.\end{solucao}
\begin{solucao}{04.45}\textbf{B.} Poder de polícia é prerrogativa juridicamente fundada; falta de competência é vício central.\end{solucao}
\begin{solucao}{04.46}\textbf{CERTO.} Motivos declarados podem vincular a validade do ato, permitindo controle de sua existência e veracidade.\end{solucao}
\begin{solucao}{04.47}\textbf{B.} A Constituição inclui pessoas jurídicas privadas prestadoras de serviço público no regime do art. 37, §6º.\end{solucao}
\begin{solucao}{04.48}\textbf{B.} Decisão que ignora argumentos relevantes pode falhar em motivação e devido processo.\end{solucao}
\begin{solucao}{04.49}\textbf{ERRADO.} Eficiência opera dentro da legalidade; não autoriza afastar exigência normativa.\end{solucao}
\begin{solucao}{04.50}\textbf{B.} A decisão de prorrogar deve ser motivada e baseada em desempenho, necessidade, vantajosidade e alternativas, não em inércia.\end{solucao}

